\documentclass[12pt]{article}
\usepackage[utf8]{inputenc}
\usepackage[T1]{fontenc}
\usepackage[margin=1in]{geometry}
\usepackage{setspace}
\usepackage{xcolor}
\usepackage[colorlinks=true, urlcolor=blue!60!black, linkcolor=blue!60!black]{hyperref}
\setlength{\parindent}{0em}
\setlength{\parskip}{0.6em}
\renewcommand{\baselinestretch}{1.0}
\begin{document}
\begin{center}
{\LARGE \textit{Tandon v. Newsom}, 593 U.S. 61 (2021)\par}
\vspace{0.5em}
{\large Case Brief --- Two-Page Edition\par}
\end{center}

\section{Memory Jogger}
California's three-household cap on in-home worship was enjoined: favoring \textit{any} comparable secular activity over religious exercise triggers strict scrutiny (\href{https://www.supremecourt.gov/opinions/20pdf/20a151_4g15.pdf}{Slip Opinion}).

\section{Detailed Case Facts}
California's COVID-19 rules capped all private home gatherings, religious or secular, at three households, while letting hair salons, retail stores, movie theaters, and indoor restaurants host more. Pastor Jeremy Wong and Karen Busch, who hosted home Bible studies for eight to twelve people, sued with other residents.

\section{Procedural History}
The district court denied a preliminary injunction, crediting experts that home gatherings were riskier. A divided Ninth Circuit denied an injunction pending appeal, comparing religious home gatherings only to secular ones. The Supreme Court granted an injunction pending appeal without argument.

\section{Judicial Votes}
5--4 \textit{per curiam} (Thomas, Alito, Gorsuch, Kavanaugh, Barrett). Roberts would have denied the application; Kagan dissented, joined by Breyer and Sotomayor.

\section{Holding}
A regulation is not neutral and generally applicable when it treats any comparable secular activity more favorably than religion; comparability depends on risk to the government's interest; the State must show that secular precautions would not suffice; and changing the rule mid-litigation does not moot the case.

\section{Analysis of Opinions}
\textbf{Per curiam:} Treating secular home gatherings equally was ``no answer''; the Ninth Circuit wrongly rejected commercial comparators and speculated about precautions, and the State had not shown that less restrictive measures would fail. \textbf{Kagan:} the right comparator is secular home gatherings, which were treated identically, and undisputed findings showed that homes are riskier because of longer interactions, poor ventilation, and less masking.

\section{Examples: Future Applications}
\textbf{Same-Side:} a police no-beard rule with medical but not religious exemptions triggers strict scrutiny. \textbf{Opposite-Side:} a noise ordinance with no secular exemptions does not. \textbf{Fence-Sitter:} a vaccine mandate allowing only medical exemptions has divided the lower courts.

\section{Critique}
\textit{Tandon} adopted Douglas Laycock's most-favored-nation theory without overruling \textit{Smith}. Critics such as Jim Oleske say it hollows out \textit{Smith} because nearly every law has secular exceptions, and Steve Vladeck faults its emergency-docket origin.

\section{Key Quotations}
\begin{quote}
``First, government regulations are not neutral and generally applicable, and therefore trigger strict scrutiny under the Free Exercise Clause, whenever they treat \textit{any} comparable secular activity more favorably than religious exercise.''
\end{quote}
\noindent --- \textit{Per Curiam}, 593 U.S. at 62.

\end{document}
